\subsection{Exact contact and antiderivative checks}
\label{contact:verification}

The fixed-coordinate correction at a negative integral power has an
elementary residue interpretation.  For $M\ge1$,
\begin{equation}\label{contact:residue}
 [x^{M-1}]\{x(K(x)-z)\}^M
 =\operatorname*{Res}_{x=0}(K(x)-z)^M\,\dd x
 =\frac{(-z+i\pi)^M-(-z-i\pi)^M}{2i\pi}.
\end{equation}
Indeed, the local coordinate $v=1/K(x)=\tan(\pi x)/\pi$ satisfies
$\dd x=\dd v/(1+\pi^2v^2)$ and has derivative $1$ at zero.  Thus
the residue is the coefficient of $v^{-1}$ in
$(v^{-1}-z)^M/(1+\pi^2v^2)$, namely
\[
 \sum_{j=0}^{\lfloor(M-1)/2\rfloor}
   \binom{M}{2j+1}(-z)^{M-2j-1}(-\pi^2)^j.
\]
The binomial theorem gives the final expression in
\eqref{contact:residue}.  Consequently the difference between the
coordinate finite-part mean and its continued spectral value is
\begin{equation}\label{contact:identity}
 \frac{(-z+i\pi)^M+(-z-i\pi)^M}{2}
       -(-z-\epsilon i\pi)^M
 =\epsilon\pi i\,[x^{M-1}]\{x(K(x)-z)\}^M.
\end{equation}
This proves the contact identity for every $M$, independently of the
general completion theorem.

The accompanying \texttt{check\_contacts.py} constructs the cotangent
Taylor jet from Bernoulli numbers and independently checks it against
the sine and cosine jets.  It then verifies \eqref{contact:identity}
exactly for $M=1,\ldots,10$ and both signs of $\epsilon$, with $z$
and $\pi$ left symbolic.  Omitting the endpoint phase produces a
nonzero polynomial residual in all twenty cases.  In addition, the
script verifies the primitive's telescoping coefficients for
$k=0,\ldots,5$ with symbolic $L\ne1$, checks the $L=1$ branch, and
reproduces the boundary example $p=2$, $u=\lambda=1$: the specialized
ordinary integrand is zero, the continued spectral value is $-1$,
and the local endpoint term has limit $-1$.  All thirty-six exact
checks and all twenty deliberate corruption controls pass.  These
finite checks complement the all-index proofs; they do not replace
them.
